\documentclass[12pt,letterpaper]{report}
\usepackage[utf8]{inputenc}
\usepackage{color}
\usepackage{fancybox}
\usepackage{fancyhdr}

\pagestyle{fancy}
\rhead{Probabilidad y Estadística}
\lhead{Tarea 01}

%\theoremstyle{definition}
%\newtheorem{defn}{Definición}[section]
%\newtheorem{conj}{Conjecture}[section]
%\newtheorem{exmp}{Example}[section]

\begin{document}
	\title{Probabilidad y Estadística: Tarea 1}
	\author{Joel C. Flores Escalante}
	\date{Agosto 2018}
	\maketitle
	
	Indicaciones: 
	
	\begin{itemize}
		\item En tu libreta de clases, rotula las hojas con los siguientes datos: tu nombre, matricula, materia, nombre de la tarea y fecha.
		\item Copia los ejercicios y preguntas en tu libreta y resuelve lo que se te pide, en tu libreta.   
		\item Toma  fotos de tu libreta con los ejercicios resueltos y pégalas en el archivo de Google Drive anexo a la tarea en classroom. Se debe observar los datos de rotulación en las fotos.
		\item IMPORTANTE: Todo el procedimiento debe estar descrito en la solución del ejercicio. No se acepta solo la respuesta.
	\end{itemize}
	
	 \begin{enumerate}
	 \item Preguntas
	 		\begin{enumerate}
	 			\item Explica que es la probabilidad
	 			\item Explica la probabilidad clásica y frencuencial
	 			\item ¿Qué es un evento y un  espacio muestral?
	 			\item Explica que es un espacio muestral discreto y uno continuo.
	 			\item Explica las operaciones básicas del diagrama de Venn: unión, intersección y complementos.
	 			\item Explica y da un ejemplo del principio multiplicativo
	 			\item Explica y da un ejemplo del principio aditivo
	 			\item Explica que es una permutación y desarrolla un ejemplo
	 			\item Explica que es una combinación y desarrolla un ejemplo
	 		\end{enumerate}

	 	\item Análisis combinatorio:
	 	\begin{enumerate}
	 		\item Cuatro libros distintos de matemáticas, seis diferentes de física y dos diferentes de química se colocan en un estante. ¿De cuántas formas distintas es posible ordenarlos si: 
	 			\begin{enumerate} 
	 				\item Los libros de cada asignatura deben estar todos juntos
	 				\item Solamente los libros de matemáticas deben estar juntos.
	 			\end{enumerate}
	 			\item De cuántas formas distintas pueden sentarse 7 personas alrededor de una mesa, si:
	 			\begin{enumerate}
	 				\item Pueden sentarse de cualquier forma.
	 				\item Dos personas determinadas no deben estar no deben estar al lado de la otra
	 			\end{enumerate}
	 			\item En una tienda deportiva tienen 7 tipos diferentes de cañas de pescar cada una en 3 colores distintos. ¿Cuantas formas de escoger una caña se tienen? 
	 			\item De un total de 5 mátematicos y 7 físicos, se forman un comité de 2 matemáticos y 3 físicos. ¿De cuántas formas puede formarse, si:
	 				\begin{enumerate}
	 					\item Si puede pertenecer a él cualquier matemático y físico.
	 					\item Un físico determinado debe pertenecer al comité
	 					\item Dos matemáticos determinados no pueden estar en el comité.
	 				\end{enumerate} 
	 			\item Con 7 consonantes y 5 vocales diferentes, ¿cuántas palabras pueden formarse, que conste de 4 consonantes y 3 vocales? No es necesario que las palabras tengan significado.	
	 			\item En cuantas formas diferentes pueden siete científicos acomodarse en una habitación triple y dos habitaciones dobles en un hotel?
	 			\item Encuéntrese el número de comités que pueden formarse con 4 químicos y 3 físicos y que comprendan 2 químicos y 1 físico.
	 			\item De cuántas maneras se pueden elegir 3 hombres de entre un grupo de 15 de forma que
	 				\begin{enumerate}
	 					\item uno de ellos debe figurar en cada grupo seleccionado
	 					\item dos de ellos no deben figurar en cada grupo seleccionado
	 					\item uno de ellos debe, y otros 2 no deben, figurar en cada grupo
	 				\end{enumerate}
	 				\item Hallar cuántos números impares de tres cifras, iguales o distintas, se pueden formar con los dígitos 3, 4, 5, 6, 7.
	 				\item En cierto sistema telefónico se utilizan cuatro letras diferentes $P,R,S,T$ y los cuatro dígitos 3, 5, 7, 8 para designar a los abonados. Hallar el máximo de <<números de teléfono>> de que puede contar dicho sistema, sabiendo que cada uno esta formado por una letra seguida de un número de cuatro cifras distintas.
	 				\item ¿Cuántos grupos se pueden formar con 8 mujeres sabiendo que en cada uno de ellos debe haber por lo menos 3?
	 				\item ¿De cuántas maneras se puede colorear un cuadro con 7 colores diferentes?
	 		\end{enumerate}
	 		
	 		Tomar en cuenta que: 
	 		\begin{itemize}
	 			\item \textit{Si una operación puede realizarse en $n_1$ formas, y si por cada una de éstas formas, entonces  las dos operaciones pueden realizarse juntas en $n_1 n_2$ formas.} 
	 			\item Permutaciones circulares: $(n-1)!$
	 			\item Número total de combinaciones de $n$ elementos $C = 2^n - 1$
	 			%\item \textit{Si un experimento puede tener cualquiera de $N$ resultados diferentes igualmente factibles, y si exactamente $n$ de estos resultados corresponden al evento $A$, entonces la probabilidad de este último es: $P(A) = \frac{n}{N}$}
	 			%\item \textit{Si $A$ y $B$ son dos eventos cualquiera, entonces $P(A \cup B) = P(A) + P(B) - P(A \cap B)$}
	 			%\item \textit{Si $A$ y $B$ son mutuamente excluyentes, entonces $ P(A \cup B) = P(A) + P(B)$}
	 		\end{itemize}
	 		
	 \end{enumerate}
	 
	
	
\end{document}